\documentclass[10pt,a4paper]{amsart}
\usepackage[margin=19mm]{geometry}
\usepackage{amsmath,amssymb,mathtools}
\usepackage[scr=rsfs,cal=euler]{mathalfa}
\usepackage{xfrac}
\usepackage[unicode,hidelinks,bookmarks=false]{hyperref}
\newcommand{\ie}{i.e.\ }
\newcommand{\cf}{cf.\ }
\newcommand{\resp}{respectively\ }
\newcommand{\Subsection}{\subsection}
\providecommand{\nobreakdash}{\nobreak\hbox{-}}
\setlength{\emergencystretch}{2em}
\multlinegap0pt
\usepackage{amssymb}
\usepackage{xcolor}
\usepackage{enumerate}
\usepackage{float}
\usepackage{tikz}
\newtheorem{thm}{Theorem}
\title{The Ptolemy-Alhazen problem with source at infinity}
\author{Masayo Fujimura, Matti Vuorinen}
\date{Source: arXiv:2603.25160v1 (CC BY 4.0)}
\begin{document}
\maketitle

\noindent\emph{Source and licence.} The Ptolemy-Alhazen problem with source at infinity. Version arXiv:2603.25160v1 (CC BY 4.0), distributed by arXiv under the Creative Commons Attribution 4.0 International licence, http://creativecommons.org/licenses/by/4.0/, as recorded in the arXiv licence metadata for this exact version and shown on the abstract page https://arxiv.org/abs/2603.25160v1. Full licence notice: \texttt{licenses/arxiv-2603.25160-CC-BY-4.0.txt}.\par\smallskip

\noindent\textit{Editorial scope.} Counted unit: the article's thm claim2 (source0.tex lines 427--433). The article's complete original proof of it is delivered with it (source0.tex lines 436--544, one \texttt{proof} environment of 109 lines). No mathematics is written, repaired or re-derived: the statement and the proof are the article's own lines, in the article's own notation, and no retained line is substituted or re-flowed.  No further source-local context is needed: every cross-reference of every retained line resolves inside the two delivered spans.  Nothing was omitted from any retained span, so the package carries no omission record.  The retained lines cite 1 works, whose entries are transcribed verbatim from the article's own bibliography source in the e-print and delivered with short editorial labels recorded in the item's reference mapping.  No retained line contains a figure environment, an image-inclusion command or drawing code, so no image is delivered and the package declares no asset dependency.  Editorial formatting only, in addition: an AMS \texttt{amsart} preamble that restates the article's own declarations and macros used by the retained lines, namely the environments \texttt{thm} on separate counters, together with the standard packages those lines need.  The retained environments print in this document as Theorem 1 in order of appearance, whereas the article prints the same items with the numbering of its own full document.  The complete native source of the article as distributed in the arXiv e-print for this exact version is bundled unchanged as \texttt{sources/197191/source0.tex}.
\begin{thm}\label{claim2}
For $ r>1 $, the four roots of the equation 
\begin{equation}\label{eq:inf}
  re^{-i\theta}w^4-w^3+w-re^{i\theta}=0
\end{equation}
lie on the unit circle.
\end{thm}

\noindent \begin{proof}
If \eqref{eq:inf} has $ w=1 $ as its solution, 
then $ \theta=0 $ follows from $ r(e^{-i\theta}-e^{i\theta})=0 $.
In this case, the corresponding parabola degenerates to a 
half-line on the real-axis because the light ray arrives 
directly above the observer.
So from now on we will assume that \eqref{eq:inf}
does not have $w=1$ as its solution.

Consider the following two equations
\begin{equation}\label{eq:D}
  (w-w_1)(w-w_2)(w-w_3)(w-w_4)=0
\end{equation}
and
\begin{equation}\label{eq:R}
  (z-z_1)(z-z_2)(z-z_3)(z-z_4)=0,
\end{equation}
where $ z_k=i\dfrac{1+w_k}{1-w_k} $ ($ k=1,2,3,4 $).
We remark that, for $ k=1,\cdots,4 $,  $ w_k\in\partial \mathbb{B}^2 $
and $ z_k\in\mathbb{R} $ are equivalent.

Suppose that $ w_1,\cdots,w_4 $ are the solutions to \eqref{eq:inf}.
Then, from Vieta's formulas, the following equalities hold
\begin{align} \notag
  &  (w_1+w_2+w_3+w_4)r-e^{i\theta}=0,\\  \notag
  &  w_1(w_2+w_3+w_4)+w_2(w_3+w_4)+w_3w_4=0, \\ \notag
  &  (w_1w_2w_3+w_1w_2w_4+w_1w_3w_4+w_2w_3w_4)r+e^{i\theta}=0, \\ \label{eq:s}
  & w_1w_2w_3w_4+e^{2i\theta}=0.
\end{align}
The equation \eqref{eq:R} can be written as 
\begin{align}\notag
  &(w_1-1)(w_2-1)(w_3-1)(w_4-1)z^4 \\ \notag
  & \quad +2i\big((w_1+w_2+w_3+w_4)-(w_1w_2w_3+w_1w_2w_4+w_1w_3w_4+w_2w_3w_4)\\
     \notag
  & \qquad
        +2w_1w_2w_3w_4-2\big)z^3 \\ \notag
  & \quad +2\big((w_1(w_2+w_3+w_4)+w_2(w_3+w_4)+w_3w_4)-3w_1w_2w_3w_4-3\big)z^2 \\
       \notag
  & \quad +2i\big((w_1+w_2+w_3+w_4)+(w_1w_2w_3+w_1w_2w_4+w_1w_3w_4+w_2w_3w_4)\\
    \notag
  & \qquad
           -2w_1w_2w_3w_4+2\big)z \\  \label{eq:R2}
  & \quad +(w_1+1)(w_2+1)(w_3+1)(w_4+1)=0
% (w1-1)*(w2-1)*(w3-1)*(w4-1)*w^4
% +2*i*((w1+w2+w3+w4)-(w1*w2*w3+w1*w2*w4+w1*w3*w4+w2*w3*w4)+2*w1*w2*w3*w4-2)*w^3
% +2*i^2*(-(w1*(w2+w3+w4)+w2*(w3+w4)+w3*w4)+3*w1*w2*w3*w4+3)*w^2
% +2*i^3*(-(w1+w2+w3+w4)+(w1*w2*w3+w1*w2*w4+w1*w3*w4+w2*w3*w4)
%   +2*w1*w2*w3*w4-2)*w +(w1+1)*(w2+1)*(w3+1)*(w4+1)*i^4
\end{align}
and the following hold,
\begin{align*}
 & (w_1-1)(w_2-1)(w_3-1)(w_4-1) \\
 & \quad =1-(w_1+w_2+w_3+w_4)
    +(w_1(w_2+w_3+w_4)+w_2(w_3+w_4)+w_3w_4) \\
 & \qquad   -(w_1w_2w_3+w_1w_2w_4+w_1w_3w_4+w_2w_3w_4)+w_1w_2w_3w_4, \\
 & (w_1+1)(w_2+1)(w_3+1)(w_4+1) \\ 
 & \quad  =1+(w_1+w_2+w_3+w_4)
      +(w_1(w_2+w_3+w_4)+w_2(w_3+w_4)+w_3w_4) \\
 & \qquad   +(w_1w_2w_3+w_1w_2w_4+w_1w_3w_4+w_2w_3w_4)+w_1w_2w_3w_4.
\end{align*}
Substituting \eqref{eq:s} into \eqref{eq:R2}, we have
\begin{align*}
  & (e^{2i\theta}-1)rz^4+4i((e^{2i\theta}+1)r-e^{i\theta})z^3-6(e^{2i\theta}-1)rz^2 \\
  & \qquad
     -4i((e^{2i\theta}+1)r+e^{i\theta})z +(e^{2i\theta}-1)r=0.
\end{align*}
Dividing the above equation by $ 2ie^{i\theta} $ gives,
\begin{equation}\label{eq:R3}
  r\sin \theta z^4+2(2r\cos \theta -1)z^3-6r\sin \theta z^2
    -2(2r\cos\theta -1)z+r\sin\theta =0.
\end{equation}

Note that equation \eqref{eq:R3} is an equation of degree four 
with real coefficients. 
It is known that the equation 
$$  ax^4 + bx^3 + cx^2 + dx + e = 0 \quad (a,b,c,d\in\mathbb{R}) $$
has four real roots if the following conditions are satisfied
(see, for example \cite{Rees}),
$$  \Delta>0,\qquad P<0 \quad\mbox{and}\quad D<0, $$ 
where
\begin{align*}
  \Delta=&
    256e^3a^3+(-192e^2db-128e^2c^2+144ed^2c-27d^4)a^2 \\
    & +\big((144e^2c-6ed^2)b^2+(-80edc^2+18d^3c)b+16ec^4-4d^2c^3\big)a \\
    & -27e^2b^4+(18edc-4d^3)b^3+(-4ec^3+d^2c^2)b^2,\\
  P=& 8ac-3b^2, \\
  D=& 64a^3e-16a^2c^2+16ab^2c-16a^2db-3b^4.
\end{align*}

For $ r>1 $, $ a=r\sin\theta$, $ b=2(2r\cos \theta -1)$, $ c=-6r\sin \theta$,
$d=-2(2r\cos\theta -1)$, and $ e=r\sin\theta$, we have
\begin{align*}
  \Delta & =  256((64\cos^6\theta+192\sin^2\theta\cos^4\theta
          +192\sin^4\theta\cos^2\theta+64\sin^6\theta)r^6\\
      & \quad +(-48\cos^4\theta-96\sin^2\theta\cos^2\theta-48\sin^4\theta)r^4
         +(12\cos^2\theta-15\sin^2\theta)r^2-1)\\
      & = 256(64r^6-48r^4+(-27\sin^2\theta+12)r^2-1)\\
      & >  256(64r^6-48r^4-15r^2-1)>256(64r^6-64r^4)>0, \\
  P   & = -12((4\cos^2\theta+4\sin^2\theta)r^2-4r\cos \theta +1)\\
      & =-12(4r^2-4r\cos \theta +1)
          <-12\big(4r(r-\cos\theta)+1\big)<0,\\
  D   & =-16((48\cos^4\theta+80\sin^2\theta\cos^2\theta+32\sin^4\theta)r^4
        +(-96\cos^3\theta-96\sin^2\theta\cos \theta)r^3 \\
      &\quad  +(72\cos^2\theta+28\sin^2\theta)r^2-24\cos \theta r+3)\\
      & = -16\bigg((4r^2+11)\Big(2r\cos \theta-\dfrac{6(4r^2+1)}{4r^2+11}\Big)^2
             +\dfrac{(4r^2-1)^2(8r^2-3)}{4r^2+11}\bigg) <0.
\end{align*}
Hence, \eqref{eq:R3} has four real roots, and the assertion is obtained.
\end{proof}

\begin{thebibliography}{99}
\bibitem[Rees]{Rees} {\sc E.~L.~Rees,} {Graphical discussion of the roots of a quartic equation}. {\em {Amer. Math. Month.}}, 29(2): 51--55, 1922. %%=========================================================================================
\end{thebibliography}
\end{document}
